% =============================================================================
%  Notas de clase — Sesión 15: Paso por Referencia
%  Programación Avanzada — Universidad Panamericana
%  Compilar: pdflatex sesion15_paso_por_referencia.tex
% =============================================================================
\documentclass[12pt, oneside]{article}
\usepackage{progavanzada}

\begin{document}

\portada{Sesión 15}{Paso por Referencia}{2026}

\tableofcontents
\newpage

% =============================================================================
\section{La trampa invisible}
% =============================================================================

Escribe esta función y dime qué esperas que imprima:

\begin{codigo}
void sumaUno(int n) {
    n = n + 1;
}

int a = 10;
sumaUno(a);
cout << a;   // ?10 u 11?
\end{codigo}

Si respondiste 11, estás en buena compañía — y en el error. El programa imprime \textbf{10}.
La variable \cpp{a} no cambió. La función modificó una copia temporal y esa copia desapareció
al terminar la llamada.

Esta sesión responde exactamente esa pregunta: ¿cómo le digo a una función que sí quiero
que modifique el original?

\begin{tecnico}[Historia: FORTRAN y el bug de los constantes]
  En \textbf{FORTRAN} (1957), el primer lenguaje de alto nivel, \textit{todos} los argumentos
  se pasaban por referencia — incluso los literales numéricos. Esto llevó a un bug legendario:
  al escribir \texttt{CALL FOO(2)} y que \texttt{FOO} modificara su parámetro, el programa
  podía alterar el valor almacenado para el literal \texttt{2}. Después de la llamada,
  cualquier aparición de \texttt{2} en el programa devolvía \texttt{3}.
  Este tipo de aliasing involuntario fue una de las razones por las que C eligió el paso
  por valor como comportamiento predeterminado.
  \par\smallskip
  \textit{Referencia:} Backus, J. (1978). ``The History of FORTRAN I, II, and III''.
  \textit{ACM SIGPLAN Notices}, 13(8), 165--180.
  \url{https://doi.org/10.1145/960118.808380}
\end{tecnico}

% =============================================================================
\section{Paso por valor: el comportamiento por defecto}
% =============================================================================

Cuando llamas a una función, C++ crea una nueva variable local en el \textit{stack}
de esa función y \textbf{copia} el valor del argumento en ella. La función trabaja
con esa copia; el original permanece intacto.

\begin{center}
\begin{tabular}{p{6cm} p{6cm}}
\toprule
\textbf{En \cpp{main()}} & \textbf{En \cpp{sumaUno()}} \\
\midrule
\begin{verbatim}
  +---------+
  |   10    |  <- int a
  +---------+
  0xAA00
\end{verbatim}
&
\begin{verbatim}
  +---------+
  |   10    |  <- int n  (copia)
  +---------+
  0xBB00
\end{verbatim}
\\
\bottomrule
\end{tabular}
\end{center}

\cpp{a} y \cpp{n} son celdas de memoria distintas. Cualquier cambio en \cpp{n} no afecta a \cpp{a}.

\begin{recuerda}
  El paso por valor es seguro y predecible: la función nunca modifica
  los datos del llamador por accidente. Es la opción correcta para datos
  pequeños (\cpp{int}, \cpp{double}, \cpp{char}) que la función solo necesita \textit{leer}.
\end{recuerda}

% =============================================================================
\section{Paso por referencia}
% =============================================================================

Para que una función \textbf{modifique} el original, declara el parámetro con
\textbf{\cpp{\&}}. El símbolo indica que la función recibe un \textit{alias}
— otro nombre para la misma celda de memoria — en lugar de una copia.

\begin{codigo}[caption={Paso por referencia: el \& en el parámetro}]
void sumaUno(int& n) {   // n es alias del argumento
    n = n + 1;
}

int a = 10;
cout << "antes:  " << a << endl;   // 10
sumaUno(a);
cout << "despues: " << a << endl;  // 11
\end{codigo}

Ahora \cpp{n} y \cpp{a} son el mismo espacio de memoria:

\begin{center}
\begin{tabular}{p{5cm} p{7cm}}
\toprule
\textbf{Paso por valor} & \textbf{Paso por referencia} \\
\midrule
\begin{verbatim}
main()     sumaUno()
+------+   +------+
|  10  |   |  10  | (copia)
+------+   +------+
0xAA00     0xBB00
\end{verbatim}
&
\begin{verbatim}
  main()
  +------+
  |  10  | <-- a == n (mismo espacio)
  +------+
  0xAA00
\end{verbatim}
\\
\bottomrule
\end{tabular}
\end{center}

\begin{advertencia}
  El \cpp{\&} en una \textbf{declaración de parámetro} (\cpp{int\& n}) indica
  paso por referencia. El \cpp{\&} en una \textbf{expresión} (\cpp{\&a}) es el
  operador dirección-de, que devuelve la dirección de \cpp{a}. Son dos usos del
  mismo símbolo con significados completamente distintos. En la \textit{llamada}
  a la función no se escribe ningún \cpp{\&}.
\end{advertencia}

% =============================================================================
\section{Paso por apuntador: el mecanismo heredado de C}
% =============================================================================

Antes de que C++ introdujera las referencias, la única forma de lograr el mismo
efecto era pasando la \textbf{dirección} de la variable y trabajando con un apuntador
dentro de la función. C todavía usa este mecanismo hoy.

\begin{codigo}[caption={Paso por apuntador}]
void sumaUno(int* p) {   // p recibe una direccion
    (*p) = (*p) + 1;     // desreferenciamos para modificar el original
}

int a = 10;
sumaUno(&a);             // pasamos la direccion de a
cout << a;               // 11
\end{codigo}

El resultado es idéntico al paso por referencia, pero la sintaxis es más ruidosa:
hay que escribir \cpp{\&a} en la llamada y \cpp{*p} dentro de la función.

\begin{ejemplo}[¿Por qué \cpp{scanf} necesita \cpp{\&x}?]
  Si alguna vez te preguntaste por qué \cpp{scanf("\%d", \&x)} exige el
  \cpp{\&} y \cpp{cin >> x} no, aquí está la respuesta: \cpp{scanf} es una
  función de C (que no tiene referencias), así que necesita la dirección para
  poder modificar \cpp{x}. \cpp{cin >>} es un operador de C++ implementado con
  referencias, por lo que la dirección se pasa de forma implícita y transparente.
  El mecanismo es el mismo; solo cambia si el compilador lo hace visible o no.
\end{ejemplo}

% =============================================================================
\section{¿Cuándo usar cada mecanismo?}
% =============================================================================

\begin{center}
\begin{tabular}{p{3.2cm} p{3.5cm} p{6cm}}
  \toprule
  \textbf{Mecanismo} & \textbf{Sintaxis} & \textbf{Cuándo usarlo} \\
  \midrule
  Por valor          & \cpp{int n}       & Dato pequeño que la función solo lee \\
  Por referencia     & \cpp{int\& n}     & La función necesita modificar el original \\
  Por ref.\ constante& \cpp{const int\& n} & Dato grande que la función solo lee (evita copias) \\
  Por apuntador      & \cpp{int* p}      & Puede ser \cpp{nullptr}, o al trabajar con arreglos y memoria dinámica \\
  \bottomrule
\end{tabular}
\end{center}

\subsection{Referencia constante: leer sin copiar}

Un \cpp{std::string} o un \cpp{std::vector} pueden ocupar kilobytes o más.
Pasarlos por valor crea una copia completa en cada llamada. La solución es
\cpp{const\&}: le prometemos al compilador que no vamos a modificar el dato,
y a cambio solo pasamos la dirección (8 bytes), no una copia.

\begin{codigo}[caption={const\& para evitar copias costosas}]
// MAL: crea una copia del string completo en cada llamada
void imprimirV(string texto) {
    cout << texto << endl;
}

// BIEN: solo pasa la direccion, sin copiar nada
void imprimir(const string& texto) {
    // texto = "otro";  // error: no se puede modificar (es const)
    cout << texto << endl;
}

string saludo = "Hola, mundo";
imprimir(saludo);   // eficiente: 8 bytes pasados, no toda la cadena
\end{codigo}

\begin{consejo}
  Regla práctica: si el parámetro es de un tipo que ocupa más de 8 bytes
  y la función no necesita modificarlo, usa \cpp{const\&}. Para tipos
  primitivos (\cpp{int}, \cpp{double}, \cpp{char}) el valor copiado es tan
  pequeño como la dirección, así que el paso por valor es igual de eficiente
  y más claro.
\end{consejo}

% =============================================================================
\section{Funciones con múltiples resultados}
% =============================================================================

Una función solo puede tener un \cpp{return}. El paso por referencia permite
"devolver" varios resultados usando \textbf{parámetros de salida}.

\begin{codigo}[caption={Mínimo y máximo de un arreglo en una sola pasada}]
void minMax(const int* arr, int n, int& minVal, int& maxVal) {
    minVal = maxVal = arr[0];
    for (int i = 1; i < n; i++) {
        if (arr[i] < minVal) minVal = arr[i];
        if (arr[i] > maxVal) maxVal = arr[i];
    }
}

int datos[6] = {3, 7, 1, 9, 4, 6};
int minimo, maximo;

minMax(datos, 6, minimo, maximo);

cout << "Minimo: " << minimo << endl;  // 1
cout << "Maximo: " << maximo << endl;  // 9
\end{codigo}

\cpp{minVal} y \cpp{maxVal} son aliases de \cpp{minimo} y \cpp{maximo}. Al escribir
en ellas dentro de la función, estamos escribiendo directamente en las variables del llamador.

\begin{tecnico}[Usos reales en la industria]
  Este patrón de parámetros de salida es omnipresente:
  \begin{itemize}
    \item La función \cpp{std::minmax\_element} de la STL regresa un par de
          iteradores usando exactamente este concepto.
    \item En motores de videojuegos (Unreal Engine, Unity via C\#), las funciones
          que calculan física o geometría frecuentemente devuelven múltiples
          vectores y normales como parámetros \cpp{out}.
    \item En sistemas embebidos, donde las excepciones están deshabilitadas,
          los parámetros de salida son la forma estándar de reportar éxito/error
          junto con el resultado.
  \end{itemize}
\end{tecnico}

% =============================================================================
\section{Arreglos: siempre por referencia implícita}
% =============================================================================

Los arreglos son un caso especial: \textbf{siempre} se pasan por referencia,
aunque no escribas ningún \cpp{\&}. Al pasar un arreglo a una función, su nombre
\textit{decae} automáticamente a un apuntador al primer elemento, y la función
trabaja directamente sobre el arreglo original.

\begin{codigo}[caption={Arreglos: la función modifica el original}]
void duplicar(int* arr, int n) {
    for (int i = 0; i < n; i++)
        arr[i] *= 2;   // modifica el arreglo original
}

int nums[5] = {1, 2, 3, 4, 5};
duplicar(nums, 5);

for (int i = 0; i < 5; i++)
    cout << nums[i] << " ";   // 2 4 6 8 10
\end{codigo}

\begin{recuerda}
  Cuando pasas un arreglo a una función siempre hay que pasar también
  el \textbf{tamaño} como parámetro adicional. Al decaer a apuntador,
  la información del tamaño se pierde: dentro de la función,
  \cpp{sizeof(arr)} daría el tamaño del apuntador (8 bytes), no del arreglo.
  Esto es exactamente lo que estudiamos con el \textit{array decay} en la
  sesión anterior.
\end{recuerda}

\subsection{Arreglos bidimensionales}

Para pasar una matriz, la función necesita conocer \textbf{todas las dimensiones
excepto la primera} (que puede omitirse):

\begin{codigo}[caption={Arreglo 2D como parámetro}]
void incrementar(int mat[][3], int filas) {
    for (int i = 0; i < filas; i++)
        for (int j = 0; j < 3; j++)
            mat[i][j]++;
}

int m[3][3] = {{1,2,3},{4,5,6},{7,8,9}};
incrementar(m, 3);
// m ahora contiene {{2,3,4},{5,6,7},{8,9,10}}
\end{codigo}

La razón por la que se necesita la segunda dimensión es aritmética de apuntadores:
para calcular \cpp{mat[i][j]}, el compilador necesita saber cuántos \cpp{int}
ocupa cada fila.

% =============================================================================
\section{Aplicación: Bubble Sort}
% =============================================================================

Bubble Sort es el ejemplo perfecto para conectar los tres conceptos de esta sesión:
el \cpp{swap} requiere paso por referencia, el \cpp{bubblesort} recibe el arreglo
como apuntador, y el \cpp{comparar} recibe por valor (solo lee).

\subsection{La función \cpp{swap}}

Intercambiar dos variables es imposible con paso por valor. Necesitamos referencias:

\begin{codigo}[caption={swap: el ejemplo canónico de paso por referencia}]
void swap(int& a, int& b) {
    int temp = a;
    a = b;
    b = temp;
}

int x = 5, y = 9;
swap(x, y);
cout << x << " " << y;  // 9 5
\end{codigo}

\begin{consejo}
  La biblioteca estándar ya incluye \cpp{std::swap} (en \texttt{<algorithm>}),
  implementado exactamente así. Internamente, todos los algoritmos de ordenamiento
  de la STL (\cpp{std::sort}, \cpp{std::stable\_sort}) la usan.
\end{consejo}

\subsection{El algoritmo completo}

La idea de Bubble Sort: compara pares adyacentes e intercámbialos si están en
el orden equivocado. Repite hasta que ningún par necesite intercambio.

\begin{center}
\begin{tabular}{cll}
  \toprule
  \textbf{Pasada} & \textbf{Arreglo} & \textbf{Acción} \\
  \midrule
  Inicial & \texttt{[5, 3, 8, 1, 4]} & --- \\
  1       & \texttt{[3, 5, 1, 4, 8]} & el 8 llegó al final \\
  2       & \texttt{[3, 1, 4, 5, 8]} & el 5 llegó a su lugar \\
  3       & \texttt{[1, 3, 4, 5, 8]} & el 3 llegó a su lugar \\
  4       & \texttt{[1, 3, 4, 5, 8]} & sin cambios: terminó \\
  \bottomrule
\end{tabular}
\end{center}

El nombre viene de que los elementos grandes ``burbujean'' hacia el final en cada pasada.

\begin{codigo}[caption={Bubble Sort completo}]
bool comparar(int a, int b) {
    return a > b;   // true si hay que intercambiar (a deberia ir despues)
}

void swap(int& a, int& b) {
    int temp = a;  a = b;  b = temp;
}

void bubblesort(int* arr, int n) {
    for (int i = 0; i < n - 1; i++)
        for (int j = 0; j < n - i - 1; j++)
            if (comparar(arr[j], arr[j+1]))
                swap(arr[j], arr[j+1]);
}
\end{codigo}

\begin{tecnico}[\cpp{std::sort} y los comparadores]
  La función \cpp{std::sort} de la STL acepta un \textbf{comparador} como
  tercer argumento — exactamente igual que nuestro \cpp{comparar}. Esto te
  permite ordenar en cualquier dirección o por cualquier criterio:
  \begin{codigo}
  #include <algorithm>
  int arr[5] = {3, 1, 4, 1, 5};

  // Orden ascendente (defecto)
  sort(arr, arr+5);

  // Orden descendente: pasa una funcion comparadora
  sort(arr, arr+5, [](int a, int b){ return a > b; });
  \end{codigo}
  La \textit{lambda} \cpp{[](int a, int b)\{ return a > b; \}} es básicamente
  nuestra función \cpp{comparar} escrita en línea. C++11 la introdujo
  precisamente para evitar tener que declarar funciones auxiliares por separado.
\end{tecnico}

% =============================================================================
\section{Ejercicios}
% =============================================================================

\subsection*{Ejercicio 1 — Elevar al cubo}

Escribe una función \cpp{cubo(int\& n)} que eleve \cpp{n} a la tercera potencia,
modificando la variable original. Pruébala con los valores $\{2, 3, 4, 5\}$ e imprime
el arreglo antes y después de aplicarla a cada elemento.

\subsection*{Ejercicio 2 — Tres resultados de una sola pasada}

Escribe una función:
\begin{codigo}
void estadisticas(const int* arr, int n,
                  int& minVal, int& maxVal, int& suma);
\end{codigo}
que calcule el mínimo, el máximo y la suma del arreglo en una sola pasada.
Pruébala con \cpp{int datos[7] = \{12, 5, 8, 21, 3, 17, 9\}}.

\subsection*{Ejercicio 3 — Normalizar un arreglo}

Escribe una función \cpp{normalizar(double* arr, int n)} que divida cada
elemento entre el valor máximo del arreglo, de modo que todos los valores queden
en el rango $[0,\,1]$. Pruébala con \cpp{double v[5] = \{20.0, 50.0, 10.0, 80.0, 35.0\}}.

\noindent\textit{Pista: necesitarás encontrar el máximo primero y luego recorrer el arreglo de nuevo.
¿Podrías reutilizar tu función del ejercicio 2?}

\subsection*{Ejercicio 4 — Bubble Sort descendente}

Modifica el \cpp{comparar} del Bubble Sort para que el algoritmo ordene en orden
\textbf{descendente}. Sin cambiar \cpp{swap} ni \cpp{bubblesort}, ¿qué línea es
la única que hay que modificar? Explica por escrito por qué basta con ese cambio.

% =============================================================================
\section{Resumen}
% =============================================================================

\begin{recuerda}
\begin{itemize}
  \item Por \textbf{valor} (\cpp{int n}): la función recibe una copia; el original no cambia.
  \item Por \textbf{referencia} (\cpp{int\& n}): la función recibe un alias; sí modifica el original.
  \item Por \textbf{referencia constante} (\cpp{const int\& n}): alias de solo lectura; útil para objetos grandes.
  \item Por \textbf{apuntador} (\cpp{int* p}): la función recibe una dirección; puede ser \cpp{nullptr}.
  \item Los arreglos \textbf{siempre} decaen a apuntador al pasarse; hay que pasar el tamaño por separado.
  \item Con referencias se pueden devolver \textbf{múltiples resultados} de una función sin usar \cpp{return}.
  \item \cpp{swap} con referencias es el patrón base de todos los algoritmos de ordenamiento.
  \item \cpp{std::sort} acepta un comparador para ordenar por cualquier criterio.
\end{itemize}
\end{recuerda}

\end{document}
